%% =====================================================================
%%  T-13-02  Weakness as a Mask
%%  -------------------------------------------------------------------
%%  One idea per file. Core is mandatory. Slots in canonical order.
%%  Max 700 words. No layout commands. No filler. New source material
%%  for this entry goes in sources/<BOOK>/T-13-02.tex -- never by
%%  rewriting this file (docs/RULES.md R0.1).
%%  =====================================================================
%% @kind: topic
%% @id: T-13-02
%% @title: Weakness as a Mask
%% @subtitle: The meek cover that invites disclosure
%% @chapter: c13
%% @order: 20
%% @status: stable
%% @sources: B1
%% @updated: 2026-09-19
%% @allow-rewrite: B1 ingest 2026-09-19 -- committed scaffold replaced by the written entry (planned -> stable lifecycle)

\topic{T-13-02}{Weakness as a Mask}{The meek cover that invites disclosure}

\begin{core}
Displayed competence triggers defence; displayed weakness lowers it. The practitioner underplays knowledge, status and sharpness --- never knows much, never argues first, slightly out of place --- because opponents who believe they are dealing with someone slower talk freely, overextend, and can be taken for more money. The mask is meekness worn deliberately.
\end{core}
\begin{mechanism}
The instruction comes from a house-to-house salesman who out-earned the tradesmen who mocked him. His rule: the man who believes he knows a lot is the easiest one to out-sharpen, and a self-styled wise guy can be taken for more money. The mechanism is the target's own self-image doing the work. A person who believes they are the sharpest in the room stops checking, discloses to demonstrate superiority, and prices you as overhead. Your displayed ignorance is the grant that lets them feel safe; their felt superiority is the hatch through which everything leaks. Underestimation also moves first moves: the confident opponent reveals positions, deadlines and doubts to someone they have classified as harmless --- the same misclassification the operator profile (T-03-02) exploits from the other side.
\end{mechanism}
\begin{conditions}
Strongest in first-contact dealings --- negotiations, sales, new rooms --- where classifications are fresh and the counterpart has an ego invested in being the expert. Weakens with repeated dealings (the mask decays as data accumulates), against counterparts who read calm-and-quiet as training rather than deficiency, and anywhere your visible competence is itself the product being bought.
\end{conditions}
\begin{application}
  \item Enter rooms one notch below your actual level: competent enough to be taken seriously, not enough to be feared.
  \item Ask the obvious question and let them teach you. Teaching feels like winning; it is disclosure.
  \item Never correct small errors. Keep your corrections for the moment they are worth a position.
  \item When they show off, take notes visibly. The notebook is a backstage pass.
  \item Time the reveal. The mask's value is realised in one moment --- the bid, the counter, the clause --- and only once.
\end{application}
\begin{feedback}
  \item Landing: they explain things you never asked about, and enjoy explaining them.
  \item Landing: terms improve or tighten in your favour while they relax --- the underestimation is pricing your offers.
  \item Failing: they start asking you technical questions. The classification is being re-checked.
  \item Failing: irony at your expense sharpens --- they suspect the mask and are probing it.
\end{feedback}
\begin{failure}
Worn among peers, the mask costs real standing: people promote and partner with competence they can see, and meekness reads as absence. It also invites genuine predation --- the display works on honest and dishonest eyes alike, and a truly predatory reader will simply attack what looks weak rather than underestimate it. And the mask can set: play dumb long enough in one room and your recorded self is dumb, whatever you reveal later.
\end{failure}
\begin{countermeasures}
  \item Recalibrate who you explain things to. The eager student may be filing everything you say.
  \item Notice who never shows their full hand while letting you show yours --- symmetric disclosure or nothing.
  \item Price people by decisions, not by presentation; the meek cover is expensive to see through and worth the effort exactly once.
  \item Audit your own rooms: if you always feel like the smartest person present, ask who is paying for that feeling (T-06-03).
\end{countermeasures}

\sources{B1}
\seesources
\seealso{T-03-02, T-07-02, T-09-03}
